\chapter{Metodologia}

\section{Materiais}

\subsection{Placa de desenvolvimento}

A placa de desenvolvimento escolhida para este trabalho foi a STM32H735G-DK, da STMicroelectronics. Essa é uma placa de demonstração da família STM32H7 que possui um display LCD-TFT integrado, projetada para facilitar o desenvolvimento e a demonstração de interfaces gráficas. Uma imagem da placa é mostrada na Figura~\ref{fig:stm32h735g-dk}.

\begin{figure}
    \centering
    \includegraphics[width=0.7\textwidth]{stm32h735g-dk.png}
    \caption{Placa de desenvolvimento STM32H735G-DK.}
    \label{fig:stm32h735g-dk}
\end{figure}

A placa é equipada com o microcontrolador STM32H735IGK6 da linha STM32H7, que possui um núcleo ARM Cortex-M7 operando a até 480 MHz, 1 MB de memória flash e 564 KB de RAM. O display integrado é um LCD-TFT de 4,3 polegadas com resolução de 480 x 272 pixels. A placa também inclui diversos periféricos e um depurador integrado para facilitar o desenvolvimento e a depuração de aplicações.

A escolha de uma placa fornecida pela STMicroelectronics se deve a qualidade das ferramentas fornecidas pela fabricante, com ampla documentação e suporte da comunidade. Bem como a possibilidade de utilizar a biblioteca gráfica TouchGFX, que pela licença só pode ser utilizada com microcontroladores STM32. A escolha do microcontrolador STM32H7 se deve ao seu alto desempenho, que permitirá avaliar o desempenho das bibliotecas gráficas sem que o hardware seja um gargalo. 

\subsection{Drivers}


\section{Métodos}

